\subsection{Inputs, Parameterization, and Cost Functions}

This study synthesizes data from academic literature, public government statistics, and operational key-informant interviews to parameterize the environment with high ecological validity \citep{McCulloch2022}.

\subsubsection{Behavioral Initialization}
The foundational psychological parameters for the household agents were derived through a meta-synthesis of existing pro-environmental behavioral studies \citep{Taraghi2025, Moeini2023}. To ground these parameters locally, initialization variables were calibrated using empirical Knowledge, Attitude, and Practices (KAP) data from riverside barangays in Northern Mindanao \citep{Paigalan2025}, which serve as a robust proxy for the coastal Municipality of Bacolod. 

To replicate the ``Intention-Action Gap'' frequently observed in developing economies \citep{Yazawa2025, Geiger2021}, agents were initialized with a relatively high Attitude ($A_0 \approx 0.66$) but a significantly lower Baseline Compliance ($B_0 \approx 0.12$ to $0.15$, scaled locally). This deliberate dissonance compels the RL agent to discover policies capable of bridging the gap between theoretical awareness and physical action.

\begin{table}[H]
    \centering
    \captionsetup{font={bf, small}} 
    \caption{Initialization Parameters for Household Agents \citep{Paigalan2025}}
    \label{tab:initialization_parameters}
    \small
    \renewcommand{\arraystretch}{1.15} 
    \begin{tabularx}{\textwidth}{ 
        >{\raggedright\arraybackslash}p{2.8cm} 
        >{\raggedright\arraybackslash}p{3.2cm} 
        c 
        c 
        >{\raggedright\arraybackslash}X 
    }
        \toprule
        \textbf{ABM Parameter} & \textbf{Source Variable} & \makecell{\textbf{Mean} \\ (1-5)} & \makecell{\textbf{Norm.} \\ (0-1)} & \textbf{Contextual Justification} \\
        \midrule
        Agent Knowledge ($K_0$) & Awareness of Programs & 3.27 & 0.65 & Moderate awareness of SWM programs but lacking technical depth. \\
        Agent Attitude ($A_0$) & Attitude on Labeling & 3.32 & 0.66 & Generally positive disposition toward segregation rules. \\
        Training Sensitivity ($S_t$) & Need for Training & 3.41 & 0.68 & Highly responsive to IEC interventions due to expressed need for training. \\
        \bottomrule
    \end{tabularx}
\end{table}

\subsubsection{Socio-demographic and Institutional Constraints}
The simulated LGU operates under strict fiscal constraints derived directly from the Municipal Environment and Natural Resources Office (MENRO). The municipality is allocated a fixed annual Solid Waste Management budget of \textpeso 1,500,000, which is discretized into a quarterly operating budget ($B_q$) of \textpeso 375,000. 

\subsubsection{Policy Scenarios and Resource Cost Functions}
The LGU's action space manipulates three primary policy levers, each carrying specific financial burdens based on regional realities \citep{Ibanez2022}:

\textbf{1. Information, Education, and Communication (IEC):} IEC costs are bifurcated into two administrative tiers. At the BLGU level (Micro-IEC), budgets fund grassroots \textit{Recoridas} ($C_{rec} = \text{\textpeso} 200$) and \textit{Purok} meetings ($C_{pur} = \text{\textpeso} 1,000$). At the municipal LGU level (Macro-IEC), the budget funds broad-reach radio broadcasts ($C_{rad} = \text{\textpeso} 500$) and large-scale assembly events ($C_{evt} = \text{\textpeso} 5,000$). The intensity of municipal IEC ($I_{macro}$) is bounded by target saturation levels:
\begin{align}
    n_{rad} &= \min \left( T_{rad}, \left\lfloor \frac{B_{lgu\_iec}}{C_{rad}} \right\rfloor \right) \label{eq:macro_radio} \\
    rem_{lgu} &= \max \left( 0, B_{lgu\_iec} - (n_{rad} \times C_{rad}) \right) \label{eq:macro_rem} \\
    n_{evt} &= \min \left( T_{evt}, \left\lfloor \frac{rem_{lgu}}{C_{evt}} \right\rfloor \right) \label{eq:macro_event} \\
    I_{macro} &= \min \left( 1.0, \omega_{rad} \left( \frac{n_{rad}}{T_{rad}} \right) + \omega_{evt} \left( \frac{n_{evt}}{T_{evt}} \right) \right) \label{eq:macro_intensity}
\end{align}

\textbf{2. Cost of Enforcement:} Modeled as a strict ``No Segregation, No Collection'' mandate \citep{Wang2023}, enforcement requires the deployment of human capital. Assuming an active quarterly working period of 66 days and the Region X minimum daily wage ($W_{Daily} = \text{\textpeso}400$), the cost is linearly tied to the number of deployed inspectors:
\begin{equation}
    C_{\mathrm{Enf}} = (N_{\mathrm{Enforcers}} \times W_{\mathrm{Daily}} \times 66)
    \label{eq:cost_enforcement}
\end{equation}

\textbf{3. Cost of Incentives:} Positive reinforcement is implemented as direct monetary or material subsidies \citep{Chen2023, Vorobeva2022}. This lever introduces a "Victim of Success" risk; as behavioral compliance increases across the grid, the municipality's financial liability increases proportionally:
\begin{equation}
    C_{\mathrm{Inc}} = (V_{\mathrm{Reward}} \times N_{\mathrm{Compliant}})
    \label{eq:cost_incentives}
\end{equation}